\DocumentMetadata{
  pdfversion=2.0,
  pdfstandard=ua-2,
  lang=en-US,
  tagging=on,
  tagging-setup={math/setup=mathml-SE,tabsorder=structure}
}
\documentclass[11pt,letterpaper]{article}
\usepackage[margin=0.68in,includefoot]{geometry}
\usepackage{newcomputermodern}
\usepackage{amsmath,amscd,mathtools}
\providecommand{\square}{\mdlgwhtsquare}
\usepackage[hidelinks]{hyperref}
\hypersetup{
  pdftitle={Numerical Linear Algebra First-Year Exam, January 2020},
  pdfauthor={University of Florida Department of Mathematics},
  pdfsubject={Graduate mathematics examination},
  pdfdisplaydoctitle=true,
  bookmarksopen=true
}
\setcounter{secnumdepth}{0}
\setlength{\parindent}{0pt}
\setlength{\parskip}{0.28em}
\setlength{\emergencystretch}{3em}
\setlength{\leftmargini}{2.15em}
\setlength{\leftmarginii}{2.65em}
\setlength{\leftmarginiii}{3em}
\renewcommand{\labelenumi}{\textbf{\theenumi.}}
\pagestyle{empty}
\raggedbottom
\allowdisplaybreaks
\newenvironment{examproblems}[1]{%
  \begin{enumerate}%
  \setcounter{enumi}{#1}%
  \setlength{\topsep}{0.35em}%
  \setlength{\partopsep}{0pt}%
  \setlength{\itemsep}{0.48em}%
  \setlength{\parsep}{0.18em}%
}{\end{enumerate}}
\newenvironment{examsubparts}{%
  \begin{enumerate}%
  \setlength{\topsep}{0.18em}%
  \setlength{\partopsep}{0pt}%
  \setlength{\itemsep}{0.16em}%
  \setlength{\parsep}{0.1em}%
}{\end{enumerate}}
\newenvironment{examinstructions}{%
  \par\begingroup\itshape
}{%
  \par\endgroup\vspace{0.18em}
}
\makeatletter
\renewcommand\section{\@startsection{section}{1}{0pt}%
  {-1.25ex plus -.25ex minus -.1ex}%
  {0.55ex plus .1ex}%
  {\normalfont\large\bfseries}}
\renewcommand{\@maketitle}{%
  \newpage\null\vskip -1.2em%
  {\footnotesize\bfseries
    \href{https://www.ufl.edu/}{University of Florida}, \href{https://math.ufl.edu/}{Department of Mathematics}\par}%
  \vskip 0.18em%
  \tagstructbegin{tag=Title}%
    {\LARGE\bfseries\href{https://gma.math.ufl.edu/exams/numerical-linear-algebra-first-year/}{Numerical Linear Algebra First-Year Exam}, January 2020\par}%
  \tagstructend%
  \vskip 0.35em%
  \tagmcbegin{artifact}%
    \hrule height 1pt%
  \tagmcend%
  \vskip 0.55em%
}
\makeatother
\title{Numerical Linear Algebra First-Year Exam, January 2020}
\author{}
\date{}
\begin{document}
\maketitle
\thispagestyle{empty}
\begin{examinstructions}
Do 4 (four) problems.
\end{examinstructions}
\begin{examproblems}{0}
\item
Let \(A \in \mathbb{C}^{m \times n}\).
\begin{examsubparts}
\item[\textnormal{(a)}]
Determine constants~\(\alpha\) and~\(\beta\) such that the following inequality holds for the~\(p\) and~\(\infty\) norms of matrix~\(A\), for integers \(p \ge 1\). Justify your answer. 
\[
\alpha \lVert A \rVert_\infty \le \lVert A \rVert_p \le \beta \lVert A \rVert_\infty.
\]
\item[\textnormal{(b)}]
Prove or give a counterexample: \(\lVert A \rVert_2 \le \lVert A \rVert_F\), where \(\lVert A \rVert_F\) is the Frobenius norm of~\(A\). If you prove this, make sure to justify each nontrivial step.
\end{examsubparts}
\item
Suppose~\(A\) is Hermitian positive definite.
\begin{examsubparts}
\item[\textnormal{(a)}]
Prove that each principal submatrix of~\(A\) is Hermitian positive definite.
\item[\textnormal{(b)}]
Prove that an element of~\(A\) with largest magnitude lies on the diagonal.
\item[\textnormal{(c)}]
Prove that~\(A\) has a Cholesky decomposition.
\end{examsubparts}
\item
Define the matrices~\(A\) and~\(B\) by 
\[
A = \begin{pmatrix}
1 & 2 & 0 \\
1 & 2 & 0 \\
1 & 2 & 0 \\
1 & 2 & 0
\end{pmatrix},
\qquad
B = \begin{pmatrix}
1 & 0 & 1 \\
0 & 2 & 0 \\
1 & 0 & 0 \\
0 & 0 & 1
\end{pmatrix}.
\]
\begin{examsubparts}
\item[\textnormal{(a)}]
Find both full and economy singular value decompositions of~\(A\).
\item[\textnormal{(b)}]
Find both full and economy QR decompositions of~\(B\).
\end{examsubparts}
\item
Let \(\lVert \cdot \rVert\) be a subordinate (induced) matrix norm.
\begin{examsubparts}
\item[\textnormal{(a)}]
If~\(E\) is \(n \times n\) with \(\lVert E \rVert < 1\), then show \(I+E\) is nonsingular and 
\[
\lVert (I+E)^{-1} \rVert \le \frac{1}{1-\lVert E \rVert}.
\]
\item[\textnormal{(b)}]
If~\(A\) is \(n \times n\) invertible and~\(E\) is \(n \times n\) with \(\lVert A^{-1} \rVert \lVert E \rVert < 1\), then show \(A+E\) is nonsingular and 
\[
\lVert (A+E)^{-1} \rVert \le \frac{\lVert A^{-1} \rVert}{1-\lVert A^{-1} \rVert \lVert E \rVert}.
\]
\end{examsubparts}
\item
Suppose the linear equation \(Ax=b\), with 
\[
A = \begin{pmatrix}
\delta & 1 \\
1 & 1
\end{pmatrix},
\qquad |\delta| < \epsilon_m/4,
\]
 is solved on a floating-point system with machine-epsilon \(\epsilon_m\), using an LU factorization (no pivoting!) followed by forward and back substitution. (You may assume the operations \((+, -, \div, \times)\), do not incur any additional errors).
\begin{examsubparts}
\item[\textnormal{(a)}]
If \(b=(1,0)^T\), compute the backward error, \(\lVert \hat{b}-b \rVert_2/\lVert b \rVert_2\), where \(\hat{b}\) is the data that satisfies \(A\hat{x}=\hat{b}\), and \(\hat{x}\) is the computed solution.
\item[\textnormal{(b)}]
Is the result of (a) sufficient to draw any conclusions about the backward-stability of the algorithm used to compute \(\hat{x}\)? Explain.
\end{examsubparts}
\end{examproblems}
\end{document}
